\chapter{ Origin of Quantum Mechanics}
\label{oqm}
{
    \setlength\epigraphwidth{0.6\textwidth}
\epigraph{ Quantum mechanics is a
different kind of a theory to represent the world.
} {R. P. Feynman}
}
The relation between classical and quantum mechanics is often
described as one of extension: quantum mechanics generalizes
classical mechanics. Actually it is correct only for a
part (though very important one) of classical
 systems ---  systems with phase space $ \mathbb{R}^{2 n} $
 and  Hamiltonian which is quadratic function of momenta,
 and beside them some relatively exotic highly symmetric systems. 
 In that case quantum mechanics
replaces classical mechanics by altering its fundamental
mathematical structure. Classical mechanics reappears only 
as some kind of  singular limit.
On the other hand  crucial components of classical theory 
like constraints, canonical transformations, trajectories,
have no faithful quantum analogues.

Non-relativistic quantum mechanics can be presented axiomatically,
without explicit reference to a classical precursor. 
However all rigorous definitions and theorems give us no way to connect 
the mathematical constructions to the real world. 
"Of course, nature knows the quantum mechanics,
and the classical mechanics is only an approximation; so there is no mystery in
the fact that in classical mechanics there is some shadow of quantum mechanical
laws---which are truly the ones underneath. To reconstruct the original object
from the shadow is not possible in any direct way, but the shadow does help
you to remember what the object looks like" \cite{flp3}. 
Quantization is the collection of procedures by which a classical mechanical 
system is associated with a non-relativistic quantum theory. 
Naturally  quantization is neither unique nor algorithmic in general 
\footnote{ Logically we should speak of de-quantization --- transition from quantum to classical system }, 
and no single scheme applies uniformly to all classical systems.Below  we analyze some quantization schemes, 
discuss no-go theorems limiting the existence of consistent quantization maps.
So the problem of \emph{quantization}---the transition from 
classical to quantum descriptions---actually works backward.

\section{Bra-Ket Notation}
\label{bra-ket-notation}

Bra--ket notation is the standard symbolic language of quantum
mechanics. While extraordinarily powerful and conceptually elegant,
it compresses multiple layers of mathematical structure into compact
expressions. 
Its success stems from its ability to encode:
\begin{itemize}
\item Linear algebraic structure,
\item Operator theory,
\item Spectral decomposition,
\item Measurement theory,
\item Tensor products and entanglement.
\end{itemize}
However, the notation often suppresses functional-analytic and
geometric subtleties. 

It appears at Paul Dirac's book \cite{dirac-pqm}  "The Principles of Quantum Mechanics" (first edition 1930). Let's describe pros and cons:
\subsection{ Pros}
\begin{enumerate}
   \item
 Coordinate-free and basis-independent
bra--ket notation expresses quantum states abstractly:
\[|\psi\rangle \in \mathcal{H}, \qquad \langle \phi| \in \mathcal{H}^*\]
This makes the formalism:
\begin{itemize}
\item
  independent of representation (position, momentum, spin, etc.)
\item
  cleanly adaptable to changes of basis
\item
  naturally compatible with symmetry arguments
\end{itemize}
Example:
\[\psi(x) = \langle x|\psi\rangle\]
The same state appears as a function only after choosing the position
basis.
\item
 Clear dual structure
  The notation explicitly encodes:
\begin{itemize}
\item
  kets: vectors \(|\psi\rangle\)
\item
  bras: dual vectors \(\langle\psi|\)
\item
  inner product:
\end{itemize}
\[\langle \phi | \psi \rangle\]
This makes conjugation and adjoints transparent:
\[(|\psi\rangle)^\dagger = \langle \psi|\]
\item
elegant operator representation:
\[A|\psi\rangle\]
outer products generate rank-one operators:
\[|\psi\rangle\langle\phi|\]
projection operators:
\[P = |a\rangle\langle a|\]
spectral decomposition:
\[A = \sum_a a |a\rangle\langle a|\]
or (continuous spectrum)
\[A = \int a\, |a\rangle\langle a|\, da\]
This makes the spectral theorem visually intuitive.
\item
Simplifies tensor products ---
composite systems are naturally written:
\[|\psi\rangle \otimes |\phi\rangle\]
Entangled states appear transparently:
\[|\Psi\rangle = \sum_{ij} c_{ij} |i\rangle \otimes |j\rangle\]
\item
Works naturally with PVM and POVM formalism
projection-valued measure:
\[E(\Delta) = \int_\Delta |x\rangle\langle x| dx\]
POVM elements:
\[F_i = M_i^\dagger M_i\]
Bra--ket notation makes these constructions compact and readable.
\end{enumerate}

\subsection{ Cons}
\begin{enumerate}
   \item
Lack of mathematical rigor (without extra care)---
Dirac freely wrote expressions like:
\[\langle x|x' \rangle = \delta(x-x').\]
But:
\begin{itemize}
\item
  \(|x\rangle \notin L^2(\mathbb{R})\)
\item
  delta functions are distributions, not functions
\item
  continuous ``eigenvectors'' are not true Hilbert space elements.
\end{itemize}
Strictly speaking, this requires the theory of rigged Hilbert spaces:
\[\Phi \subset \mathcal{H} \subset \Phi',\]
Without this, the notation hides subtle functional analysis.
   \item
Can hide domain issues ---
unbounded operators (like momentum or Hamiltonians) require domain care:
\[P = -i\hbar \frac{d}{dx}.\]
In bra--ket notation, one may forget:
\begin{itemize}
\item
  domain restrictions
\item
  self-adjoint extensions
\item
  boundary conditions.
\end{itemize}
These are essential in rigorous analysis.
   \item
Ambiguity in infinite dimensions---
in finite dimensions:
\[\langle \psi| = |\psi\rangle^\dagger\]
in infinite dimensions:
\begin{itemize}
\item
  The identification depends on Riesz representation.
\item
  Not all linear functionals are representable without topology
  assumptions.
\end{itemize}
The notation may conceal these subtleties.
   \item
 Overuse leads to symbolic manipulation without understanding.
Some common problems:
\begin{itemize}
\item
  Treating kets like column vectors everywhere.
\item
  Confusing formal identities with rigorous equalities.
\item
  Ignoring operator domains and closure.
\end{itemize}
It encourages formal calculation over structural analysis.
\end{enumerate}

\subsection{Some Explanations}
\label{some-explanations}

Dirac's book \cite{dirac-pqm} contains
statements that are mathematically unclear, incomplete, or (by modern standards) incorrect. 
But he did not work in the modern functional-analytic framework that developed later.
For example:
\begin{itemize}
   \item He freely used “eigenvectors” of operators with continuous spectrum.
    \item He manipulated $ \delta $-functions before distribution theory existed.
    \item He treated unbounded operators without domain analysis.
    \item He wrote symbolic identities like \[ \langle x|p \rangle=\exp(i p x/\hbar) \]
as if everything were finite-dimensional.
\end{itemize}
From today’s perspective, this is not rigorous.
From 1930’s perspective, it was revolutionary.

When Dirac wrote his book:
\begin{itemize}
   \item Hilbert space theory was still developing.
\item Spectral theorem was not yet standard physics language.
\item Rigged Hilbert spaces did not exist.
\item Distribution theory did not exist.
\end{itemize}
He was not misusing established mathematics.
He was inventing a new symbolic calculus before the mathematics to justify it existed.
Later, mathematicians 
provided rigorous frameworks that made Dirac’s manipulations meaningful.

However some of statements are genuinely incorrect, for example:
\begin{itemize}
   \item implicit assumption that every observable has a complete set of eigenvectors;
\item overuse of spectral decompositions without domain care;
\item loose treatment of continuous spectra;
\item ambiguities about self-adjoint vs symmetric operators;
\item (most serious!) statement that canonical transformations in classical mechanics 
    may be lifted to unitary transformations in quantum mechanics.
\end{itemize}

In his book \cite{neumann-mfqm}  Mathematische Grundlagen der Quantenmechanik (1932) von Neumann made everything measure-theoretic and Hilbert-space rigorous and avoided $ \delta $-functions entirely.
Though von Neumann’s formulation is mathematically precise,
Dirac’s notation became universal in physics.
Because Dirac’s formalism is vastly more flexible and computationally powerful.

Today we understand that Dirac’s formalism is rigorously captured by:
\begin{itemize}
   \item rigged Hilbert spaces;
\item spectral theorem for unbounded self-adjoint operators;
\item distribution theory
\end{itemize}
and there is no reason not to use it.

\section{Dirac Quantization}
\label{dirac-quantization}
However  not all Dirac's ideas prove to be right. He sees
quantization as a map of classical Hamiltonian system into
quantum one.

A classical non-relativistic system is described by a symplectic manifold $(M,\omega)$, 
typically $M=T^*Q$ with canonical symplectic form
\begin{equation}
\omega = dq^i \wedge dp_i.
\end{equation}
Observables are smooth functions $f \in C^\infty(M)$ equipped with the Poisson bracket
\begin{equation}
{f,g} = \omega(X_f,X_g),
\end{equation}
where $X_f$ is the Hamiltonian vector field defined by $\iota_{X_f}\omega = df$. 
The resulting Poisson algebra encodes both kinematics and dynamics.

Time evolution is generated by a Hamiltonian function $H$ according to
\begin{equation}
\frac{df}{dt} = \{f,H\}.
\end{equation}
This structure motivates the canonical correspondence between Poisson brackets and 
commutators in quantum theory.

One often idealizes quantization as a map
\begin{equation}
Q : C^\infty(M) \to \mathcal{A}(\mathcal{H}),
\end{equation}
from classical observables to operators on a Hilbert space $\mathcal{H}$, satisfying 
linearity, reality, normalization, and the Poisson--commutator correspondence
\begin{equation}
Q(\{f,g\}) = \frac{1}{i\hbar}[Q(f),Q(g)].
\end{equation}

However the Groenewold--Van Hove theorem shows that no such map exists on the full 
Poisson algebra of observables for generic phase spaces. 
Even for $M=\mathbb{R}^{2n}$, the correspondence can be consistently imposed only on 
restricted subalgebras \footnote{ In my paper \cite{panfil-q} I studied these subalgebras},
demonstrating that quantization is intrinsically limited.

\section {Canonical Quantization}
\label{canonical-quantization}

Canonical quantization assigns self-adjoint operators $\hat q^i, \hat p_j$ on Hilbert space $ \mathcal{H} $
satisfying canonical commutation relations
\begin{equation}
[\hat q^i, \hat p_j] = i\hbar , \delta^i{}_j.
\end{equation}
to classical observables $q^i, p_j$ --- coordinates on phase space $ \mathbb{R}^{2 n} $.
Though this prescription is insufficient for general systems and obscures deeper structural issues: 
the relation between classical observables and quantum operators, the (non-)existence of functorial 
quantization procedures, and the origin of inequivalent quantizations. 

For most practical problems it is realized as
\begin{equation}
q^i \mapsto \hat q^i = q^i, \qquad p_j \mapsto \hat p_j = -i\hbar \frac{\partial}{\partial q^j},
\end{equation}
acting on $\mathcal{H}=L^2(Q,dq)$. For general Hamiltonians, operator ordering ambiguities arise, 
and additional analytic input is required to ensure self-adjointness and well-defined dynamics.
The standard name "canonical quantization" is misleading,
but generally used.

